\section{VisualWebArena: un banco de pruebas de evaluación orientado a tareas reales}

\subsection{De WebArena a VisualWebArena}
WebArena ya estructuró las tareas de interacción con páginas web, pero la representación puramente textual tiene limitaciones evidentes: gran parte de la semántica visual (el diseño, las imágenes, los gráficos, la posición de los botones) no se puede conservar por completo. VisualWebArena, en cambio, incorpora explícitamente la información visual y exige que el agente entienda a la vez las capturas de pantalla de la página y los elementos estructurados de la web.

\begin{importantbox}{Qué resuelve VisualWebArena}
\begin{itemize}
    \item las tareas ya no dependen solo del texto HTML, sino del contexto visual real;
    \item la evaluación ya no mide solo «si se puede escribir una respuesta textual», sino «si se puede completar una trayectoria de interacción»;
    \item el origen de los errores se acerca más al de los sistemas reales, e incluye clics equivocados, lecturas erróneas, detecciones omitidas y fallas de retroceso.
\end{itemize}
\end{importantbox}

\begin{knowledgebox}{Los tipos de tarea suelen incluir tres categorías}
\begin{enumerate}
    \item \textbf{Tipo de recuperación visual}: localizar el objetivo a partir de una imagen o un diseño;
    \item \textbf{Tipo de flujo entre sitios}: completar la tarea encadenando varias pestañas y varios sitios web;
    \item \textbf{Tipo de satisfacción de condiciones}: cumplir restricciones como precio, calificación o tiempo, y completar el envío.
\end{enumerate}
\end{knowledgebox}

\subsection{Por qué las tareas ``fáciles para los humanos'' siguen siendo difíciles para el agente}
El expositor mencionó que muchas tareas no son complejas para los humanos, pero siguen siendo difíciles para el agente; la razón principal es que ``la cadena de decisiones es demasiado larga y los errores se acumulan''. Un descuido aparentemente pequeño (hacer clic en el botón equivocado, leer mal un elemento de la página) puede desviar por completo la trayectoria posterior.

\begin{warningbox}{Las trampas de las métricas de evaluación}
Si solo se observa la tasa de éxito final, es fácil pasar por alto los modos de falla intermedios; si solo se observan los pasos intermedios, se puede sobrestimar la capacidad final de entrega. El método más confiable es una evaluación conjunta de ``métrica final + calidad de la trayectoria + tipo de error''.
\end{warningbox}

\begin{table}[H]
\centering
\small
\begin{tabular}{p{2.8cm}p{4.0cm}p{4.2cm}}
\toprule
\textbf{Dimensión de evaluación} & \textbf{Práctica común de un LLM de una sola ronda} & \textbf{Práctica al estilo VisualWebArena} \\
\midrule
Forma de entrada & prompt de texto puro & captura de pantalla de la página + estado estructurado de la web \\
Forma de salida & respuesta textual & secuencia de acciones de varios pasos \\
Criterio de corrección & coincidencia textual & resultado final de la tarea + consistencia de la trayectoria \\
Forma principal de falla & alucinación/omisión & operación errónea/falla de retroceso/desviación de la estrategia \\
\bottomrule
\end{tabular}
\caption{El cambio de paradigma de la evaluación de preguntas y respuestas textuales a la evaluación de tareas interactivas}
\end{table}

\subsection{Resumen de la sección}
El valor de VisualWebArena radica en transformar el objetivo de evaluación de ``saber decir'' a ``saber hacer'', e integrar en un marco unificado la percepción multimodal, la toma de decisiones de cadena larga y la estabilidad de ejecución.

